% chapter02.tex — Numbers of Various Sorts (source: calcu5.pdf).
% Book pp. 21–35. PDF map: calcu5 p.23 = book p.21, then sequential
% through p.37 = book p.35 (no swapped leaves in this half of the unit).

\spivakchapter{2}{Numbers of Various Sorts}

In Chapter~1 we used the word ``number'' very loosely, despite our concern
with the basic properties of numbers. It will now be necessary to distinguish
carefully various kinds of numbers.

The simplest numbers are the ``counting numbers''
\[ 1, 2, 3, \dots. \]
The fundamental significance of this collection of numbers is emphasized by
its symbol $\mathbf{N}$ (for \textbf{natural} numbers). A brief glance at
P1--P12 will show that our basic properties of ``numbers'' do not apply to
$\mathbf{N}$---for example, P2 and P3 do not make sense for $\mathbf{N}$. From
this point of view the system $\mathbf{N}$ has many deficiencies.
Nevertheless, $\mathbf{N}$ is sufficiently important to deserve several
comments before we consider larger collections of numbers.

The most basic property of $\mathbf{N}$ is the principle of ``mathematical
induction.'' Suppose $P(x)$ means that the property $P$ holds for the number
$x$. Then the principle of mathematical induction states that $P(x)$ is true
for all natural numbers $x$ provided that

\begin{steps}
\item $P(1)$ is true.
\item Whenever $P(k)$ is true, $P(k+1)$ is true.
\end{steps}

Note that condition~(2) merely asserts the truth of $P(k+1)$ under the
assumption that $P(k)$ is true; this suffices to ensure the truth of $P(x)$
for all $x$, if condition~(1) also holds. In fact, if $P(1)$ is true, then it
follows that $P(2)$ is true (by using~(2) in the special case $k = 1$). Now,
since $P(2)$ is true it follows that $P(3)$ is true (using~(2) in the special
case $k = 2$). It is clear that each number will eventually be reached by a
series of steps of this sort, so that $P(k)$ is true for all numbers $k$.

A favorite illustration of the reasoning behind mathematical induction
envisions an infinite line of people,
\[ \text{person number 1, person number 2, person number 3, } \dots. \]
If each person has been instructed to tell any secret he hears to the person
behind him (the one with the next largest number) and a secret is told to
person number~1, then clearly every person will eventually learn the secret.
If $P(x)$ is the assertion that person number $x$ will learn the secret, then
the instructions given (to tell all secrets learned to the next person)
assures that condition~(2) is true, and telling the secret to person number~1
makes~(1) true. The following example is a less facetious use of mathematical
induction. There is a useful and striking formula which expresses the sum of
the first $n$ numbers in a simple way:
\[
1 + \cdots + n = \frac{n(n+1)}{2}.
\]
To prove this formula, note first that it is clearly true for $n = 1$. Now
\emph{assume} that for some natural number $k$ we have
\[
1 + \cdots + k = \frac{k(k+1)}{2}.
\]
Then
\begin{align*}
1 + \cdots + k + (k+1)
  &= \frac{k(k+1)}{2} + k + 1\\
  &= \frac{k(k+1) + 2k + 2}{2}\\
  &= \frac{k^2 + 3k + 2}{2}\\
  &= \frac{(k+1)(k+2)}{2},
\end{align*}
so the formula is also true for $k + 1$. By the principle of induction this
proves the formula for all natural numbers $n$. This particular example
illustrates a phenomenon that frequently occurs, especially in connection
with formulas like the one just proved. Although the proof by induction is
often quite straightforward, the method by which the formula was discovered
remains a mystery. Problems~5 and~6 indicate how some formulas of this type
may be derived.

The principle of mathematical induction may be formulated in an equivalent
way without speaking of ``properties'' of a number, a term which is
sufficiently vague to be eschewed in a mathematical discussion. A more precise
formulation states that if $A$ is any collection (or ``set''---a synonymous
mathematical term) of natural numbers and

\begin{steps}
\item $1$ is in $A$,
\item $k + 1$ is in $A$ whenever $k$ is in $A$,
\end{steps}

then $A$ is the set of all natural numbers. It should be clear that this
formulation adequately replaces the less formal one given previously---we just
consider the set $A$ of natural numbers $x$ which satisfy $P(x)$. For example,
suppose $A$ is the set of natural numbers $n$ for which it is true that
\[
1 + \cdots + n = \frac{n(n+1)}{2}.
\]
Our previous proof of this formula showed that $A$ contains~$1$, and that
$k+1$ is in $A$, if $k$ is. It follows that $A$ is the set of all natural
numbers, i.e., that the formula holds for all natural numbers $n$.

There is yet another rigorous formulation of the principle of mathematical
induction, which looks quite different. If $A$ is any collection of natural
numbers, it is tempting to say that $A$ must have a smallest member.
Actually, this statement can fail to be true in a rather subtle way. A
particularly important set of natural numbers is the collection $A$ that
contains no natural numbers at all, the ``empty collection'' or ``null
set,''\footnote{Although it may not strike you as a collection, in the
ordinary sense of the word, the null set arises quite naturally in many
contexts. We frequently consider the set $A$, consisting of all $x$ satisfying
some property $P$; often we have no guarantee that $P$ is satisfied by
\emph{any} number, so that $A$ might be $\emptyset$---in fact often one proves
that $P$ is always false by showing that $A = \emptyset$.}
denoted by $\emptyset$. The null set $\emptyset$ is a collection of natural
numbers that has no smallest member---in fact, it has no members at all. This
is the only possible exception, however; if $A$ is a nonnull set of natural
numbers, then $A$ has a least member. This ``intuitively obvious'' statement,
known as the ``well-ordering principle,'' can be proved from the principle of
induction as follows. Suppose that the set $A$ has no least member. Let $B$ be
the set of natural numbers $n$ such that $1, \dots, n$ are all \emph{not} in
$A$. Clearly $1$ is in $B$ (because if $1$ were in $A$, then $A$ would have $1$
as smallest member). Moreover, if $1, \dots, k$ are not in $A$, surely $k+1$
is not in $A$ (otherwise $k+1$ would be the smallest member of $A$), so
$1, \dots, k+1$ are all not in $A$. This shows that if $k$ is in $B$, then
$k+1$ is in $B$. It follows that every number $n$ is in $B$, i.e., the numbers
$1, \dots, n$ are \emph{not} in $A$ for any natural number $n$. Thus
$A = \emptyset$, which completes the proof.

It is also possible to prove the principle of induction from the well-ordering
principle (Problem~10). Either principle may be considered as a basic
assumption about the natural numbers.

There is still another form of induction which should be mentioned. It
sometimes happens that in order to prove $P(k+1)$ we must assume not only
$P(k)$, but also $P(l)$ for all natural numbers $l \leq k$. In this case we
rely on the ``principle of complete induction'': If $A$ is a set of natural
numbers and

\begin{steps}
\item $1$ is in $A$,
\item $k + 1$ is in $A$ if $1, \dots, k$ are in $A$,
\end{steps}

then $A$ is the set of all natural numbers.

Although the principle of complete induction may appear much stronger than
the ordinary principle of induction, it is actually a consequence of that
principle. The proof of this fact is left to the reader, with a hint
(Problem~11). Applications will be found in Problems~7, 17, 20 and~22.

Closely related to proofs by induction are ``recursive definitions.'' For
example, the number $n!$ (read ``$n$ factorial'') is defined as the product of
all the natural numbers less than or equal to $n$:
\[ n! = 1 \cdot 2 \cdot \ldots \cdot (n-1) \cdot n. \]
This can be expressed more precisely as follows:
\begin{align*}
&(1) & 1! &= 1\\
&(2) & n! &= n \cdot (n-1)!.
\end{align*}
This form of the definition exhibits the relationship between $n!$ and
$(n-1)!$ in an explicit way that is ideally suited for proofs by induction.
Problem~23 reviews a definition already familiar to you, which may be
expressed more succinctly as a recursive definition; as this problem shows,
the recursive definition is really necessary for a rigorous proof of some of
the basic properties of the definition.

One definition which may not be familiar involves some convenient notation
which we will constantly be using. Instead of writing
\[ a_1 + \cdots + a_n, \]
we will usually employ the Greek letter $\Sigma$ (capital sigma, for ``sum'')
and write
\[ \sum_{i=1}^{n} a_i. \]
In other words, $\sum_{i=1}^{n} a_i$ denotes the sum of the numbers obtained
by letting $i = 1, 2, \dots, n$. Thus
\[
\sum_{i=1}^{n} i = 1 + 2 + \cdots + n = \frac{n(n+1)}{2}.
\]
Notice that the letter $i$ really has nothing to do with the number denoted by
$\sum_{i=1}^{n} i$, and can be replaced by any convenient symbol (except $n$,
of course!):
\begin{align*}
\sum_{j=1}^{n} j &= \frac{n(n+1)}{2},\\
\sum_{j=1}^{i} j &= \frac{i(i+1)}{2},\\
\sum_{n=1}^{j} n &= \frac{j(j+1)}{2}.
\end{align*}
To define $\sum_{i=1}^{n} a_i$ precisely really requires a recursive
definition:
\begin{align*}
&(1) & \sum_{i=1}^{1} a_i &= a_1,\\
&(2) & \sum_{i=1}^{n} a_i &= \sum_{i=1}^{n-1} a_i + a_n.
\end{align*}
But only purveyors of mathematical austerity would insist too strongly on
such precision. In practice, all sorts of modifications of this symbolism are
used, and no one ever considers it necessary to add any words of explanation.
The symbol
\[ \smashoperator{\sum_{\substack{i=1\\ i \neq 4}}^{n}} a_i, \]
for example, is an obvious way of writing
\[ a_1 + a_2 + a_3 + a_5 + a_6 + \cdots + a_n, \]
or more precisely,
\[ \sum_{i=1}^{3} a_i + \sum_{i=5}^{n} a_i. \]

The deficiencies of the natural numbers which we discovered at the beginning
of this chapter may be partially remedied by extending this system to the set
of \textbf{integers}
\[ \dots, -2, -1, 0, 1, 2, \dots. \]
This set is denoted by $\mathbf{Z}$ (from German ``Zahl,'' number). Of
properties P1--P12, only P7 fails for $\mathbf{Z}$.

A still larger system of numbers is obtained by taking quotients $m/n$ of
integers (with $n \neq 0$). These numbers are called \textbf{rational
numbers}, and the set of all rational numbers is denoted by $\mathbf{Q}$ (for
``quotients''). In this system of numbers all of P1--P12 are true. It is
tempting to conclude that the ``properties of numbers,'' which we studied in
some detail in Chapter~1, refer to just one set of numbers, namely,
$\mathbf{Q}$. There is, however, a still larger collection of numbers to which
properties P1--P12 apply---the set of all \textbf{real numbers}, denoted by
$\mathbf{R}$. The real numbers include not only the rational numbers, but
other numbers as well (the \textbf{irrational numbers}) which can be
represented by infinite decimals; $\pi$ and $\sqrt{2}$ are both examples of
irrational numbers. The proof that $\pi$ is irrational is not easy---we shall
devote all of Chapter~16 of Part~III to a proof of this fact. The
irrationality of $\sqrt{2}$, on the other hand, is quite simple, and was known
to the Greeks. (Since the Pythagorean theorem shows that an isosceles right
triangle, with sides of length~$1$, has a hypotenuse of length $\sqrt{2}$, it
is not surprising that the Greeks should have investigated this question.)
The proof depends on a few observations about the natural numbers. Every
natural number $n$ can be written either in the form $2k$ for some integer
$k$, or else in the form $2k + 1$ for some integer $k$ (this ``obvious'' fact
has a simple proof by induction (Problem~8)). Those natural numbers of the
form $2k$ are called \textbf{even}; those of the form $2k+1$ are called
\textbf{odd}. Note that even numbers have even squares, and odd numbers have
odd squares:
\begin{align*}
(2k)^2 &= 4k^2 = 2 \cdot (2k^2),\\
(2k+1)^2 &= 4k^2 + 4k + 1 = 2 \cdot (2k^2 + 2k) + 1.
\end{align*}
In particular it follows that the converse must also hold: if $n^2$ is even,
then $n$ is even; if $n^2$ is odd, then $n$ is odd. The proof that $\sqrt{2}$
is irrational is now quite simple. Suppose that $\sqrt{2}$ were rational; that
is, suppose there were natural numbers $p$ and $q$ such that
\[ \left(\frac{p}{q}\right)^2 = 2. \]
We can assume that $p$ and $q$ have no common divisor (since all common
divisors could be divided out to begin with). Now we have
\[ p^2 = 2q^2. \]
This shows that $p^2$ is even, and consequently $p$ must be even; that is,
$p = 2k$ for some natural number $k$. Then
\[ p^2 = 4k^2 = 2q^2, \]
so
\[ 2k^2 = q^2. \]
This shows that $q^2$ is even, and consequently that $q$ is even. Thus both
$p$ and $q$ are even, contradicting the fact that $p$ and $q$ have no common
divisor. This contradiction completes the proof.

It is important to understand precisely what this proof shows. We have
demonstrated that there is no rational number $x$ such that $x^2 = 2$. This
assertion is often expressed more briefly by saying that $\sqrt{2}$ is
irrational. Note, however, that the use of the symbol $\sqrt{2}$ implies the
existence of \emph{some} number (necessarily irrational) whose square is~$2$.
We have not proved that such a number exists and we can assert confidently
that, at present, a proof is \emph{impossible} for us. Any proof at this
stage would have to be based on P1--P12 (the only properties of $\mathbf{R}$
we have mentioned); since P1--P12 are also true for $\mathbf{Q}$ the exact
same argument would show that there is a rational number whose square is~$2$,
and this we know is false. (Note that the reverse argument will not
work---our proof that there is no rational number whose square is~$2$ cannot
be used to show that there is no real number whose square is~$2$, because our
proof used not only P1--P12 but also a special property of $\mathbf{Q}$, the
fact that every number in $\mathbf{Q}$ can be written $p/q$ for integers $p$
and $q$.)

This particular deficiency in our list of properties of the real numbers
could, of course, be corrected by adding a new property which asserts the
existence of square roots of positive numbers. Resorting to such a measure
is, however, neither aesthetically pleasing nor mathematically satisfactory;
we would still not know that every number has an $n$th root if $n$ is odd, and
that every positive number has an $n$th root if $n$ is even. Even if we
assumed this, we could not prove the existence of a number $x$ satisfying
$x^5 + x + 1 = 0$ (even though there does happen to be one), since we do not
know how to write the solution of the equation in terms of $n$th roots (in
fact, it is known that the solution cannot be written in this form). And, of
course, we certainly do not wish to assume that all equations have solutions,
since this is false (no real number $x$ satisfies $x^2 + 1 = 0$, for example).
In fact, this direction of investigation is not a fruitful one. The most
useful hints about the property distinguishing $\mathbf{R}$ from
$\mathbf{Q}$, the most compelling evidence for the necessity of elucidating
this property, do not come from the study of numbers alone. In order to study
the properties of the real numbers in a more profound way, we must study more
than the real numbers. At this point we must begin with the foundations of
calculus, in particular the fundamental concept on which calculus is
based---functions.

\subsection*{PROBLEMS}

\pnum{1.} Prove the following formulas by induction.
\begin{romans}
\item $1^2 + \cdots + n^2 = \dfrac{n(n+1)(2n+1)}{6}$.
\item $1^3 + \cdots + n^3 = (1 + \cdots + n)^2$.
\end{romans}

\pnum{2.} Find a formula for
\begin{romans}
\item $\displaystyle\sum_{i=1}^{n} (2i - 1) = 1 + 3 + 5 + \cdots + (2n - 1)$.
\item $\displaystyle\sum_{i=1}^{n} (2i - 1)^2 = 1^2 + 3^2 + 5^2 + \cdots + (2n - 1)^2$.
\end{romans}
\textit{Hint:} What do these expressions have to do with
$1 + 2 + 3 + \cdots + 2n$ and $1^2 + 2^2 + 3^2 + \cdots + (2n)^2$?

\pnum{3.} If $0 \leq k \leq n$, the ``binomial coefficient'' $\dbinom{n}{k}$ is
defined by
\[
\binom{n}{k} = \frac{n!}{k!(n-k)!}
= \frac{n(n-1)\cdots(n-k+1)}{k!},
\qquad \text{if } k \neq 0, n
\]
\[
\binom{n}{0} = \binom{n}{n} = 1
\]
(a special case of the first formula if we define $0! = 1$),
and for $k < 0$ or $k > n$ we just define the binomial coefficient to be~$0$.

\begin{letters}
\item Prove that
\[ \binom{n+1}{k} = \binom{n}{k-1} + \binom{n}{k}. \]
(The proof does not require an induction argument.)

This relation gives rise to the following configuration, known as
``Pascal's triangle''---a number not on one of the sides is the sum of the
two numbers above it; the binomial coefficient $\binom{n}{k}$ is the
$(k+1)$st number in the $(n+1)$st row.
\begin{center}
\begin{tabular}{ccccccccccc}
&&&&& $1$\\
&&&& $1$ && $1$\\
&&& $1$ && $2$ && $1$\\
&& $1$ && $3$ && $3$ && $1$\\
& $1$ && $4$ && $6$ && $4$ && $1$\\
$1$ && $5$ && $10$ && $10$ && $5$ && $1$
\end{tabular}
\end{center}

\item Notice that all the numbers in Pascal's triangle are natural numbers.
  Use part~(a) to prove by induction that $\binom{n}{k}$ is always a natural
  number. (Your entire proof by induction will, in a sense, be summed up in a
  glance by Pascal's triangle.)
\item Give another proof that $\binom{n}{k}$ is a natural number by showing
  that $\binom{n}{k}$ is the number of sets of exactly $k$ integers each
  chosen from $1, \dots, n$.
\item Prove the ``binomial theorem'': If $a$ and $b$ are any numbers and $n$
  is a natural number, then
\begin{align*}
(a+b)^n
  &= a^n + \binom{n}{1} a^{n-1} b + \binom{n}{2} a^{n-2} b^2 + \cdots
     + \binom{n}{n-1} a b^{n-1} + b^n \\
  &= \sum_{j=0}^{n} \binom{n}{j} a^{n-j} b^j.
\end{align*}
\item Prove that
\begin{romans}
\item $\displaystyle\sum_{j=0}^{n} \binom{n}{j}
      = \binom{n}{0} + \cdots + \binom{n}{n} = 2^n$.
\item $\displaystyle\sum_{j=0}^{n} (-1)^j \binom{n}{j}
      = \binom{n}{0} - \binom{n}{1} + \cdots \pm \binom{n}{n} = 0$.
\item $\displaystyle\sum_{l \text{ odd}} \binom{n}{l}
      = \binom{n}{1} + \binom{n}{3} + \cdots = 2^{n-1}$.
\item $\displaystyle\sum_{l \text{ even}} \binom{n}{l}
      = \binom{n}{0} + \binom{n}{2} + \cdots = 2^{n-1}$.
\end{romans}
\end{letters}

\pnum{4.}
\begin{letters}
\item Prove that
\[
\sum_{k=0}^{l} \binom{n}{k} \binom{m}{l-k} = \binom{n+m}{l}.
\]
\textit{Hint:} Apply the binomial theorem to $(1+x)^n(1+x)^m$.
\item Prove that
\[
\sum_{k=0}^{n} \binom{n}{k}^2 = \binom{2n}{n}.
\]
\end{letters}

\pnum{5.}
\begin{letters}
\item Prove by induction on $n$ that
\[ 1 + r + r^2 + \cdots + r^n = \frac{1 - r^{n+1}}{1 - r} \]
if $r \neq 1$ (if $r = 1$, evaluating the sum certainly presents no problem).
\item Derive this result by setting $S = 1 + r + \cdots + r^n$, multiplying
  this equation by $r$, and solving the two equations for $S$.
\end{letters}

\pnum{6.} The formula for $1^2 + \cdots + n^2$ may be derived as follows. We
begin with the formula
\[ (k+1)^3 - k^3 = 3k^2 + 3k + 1. \]
Writing this formula for $k = 1, \dots, n$ and adding, we obtain
\[
\begin{array}{rcl}
2^3 - 1^3 &=& 3 \cdot 1^2 + 3 \cdot 1 + 1\\
3^3 - 2^3 &=& 3 \cdot 2^2 + 3 \cdot 2 + 1\\
&\vdots\\
(n+1)^3 - n^3 &=& 3 \cdot n^2 + 3 \cdot n + 1\\
\hline
(n+1)^3 - 1 &=& 3[1^2 + \cdots + n^2] + 3[1 + \cdots + n] + n.
\end{array}
\]
Thus we can find $\sum_{k=1}^{n} k^2$ if we already know $\sum_{k=1}^{n} k$
(which could have been found in a similar way). Use this method to find
\begin{romans}
\item $1^3 + \cdots + n^3$.
\item $1^4 + \cdots + n^4$.
\item $\dfrac{1}{1\cdot 2} + \dfrac{1}{2\cdot 3} + \cdots + \dfrac{1}{n(n+1)}$.
\item $\dfrac{3}{1^2 \cdot 2^2} + \dfrac{5}{2^2 \cdot 3^2} + \cdots +
  \dfrac{2n+1}{n^2(n+1)^2}$.
\end{romans}

\pnum{*7.} Use the method of Problem~6 to show that $\sum_{i=1}^{n} k^p$ can
always be written in the form
\[ \frac{n^{p+1}}{p+1} + A n^p + B n^{p-1} + C n^{p-2} + \cdots. \]
(The first 10 such expressions are
\begin{align*}
\sum_{k=1}^{n} k
  &= \tfrac12 n^2 + \tfrac12 n\\
\sum_{k=1}^{n} k^2
  &= \tfrac13 n^3 + \tfrac12 n^2 + \tfrac16 n\\
\sum_{k=1}^{n} k^3
  &= \tfrac14 n^4 + \tfrac12 n^3 + \tfrac14 n^2\\
\sum_{k=1}^{n} k^4
  &= \tfrac15 n^5 + \tfrac12 n^4 + \tfrac13 n^3 - \tfrac1{30} n\\
\sum_{k=1}^{n} k^5
  &= \tfrac16 n^6 + \tfrac12 n^5 + \tfrac5{12} n^4 - \tfrac1{12} n^2\\
\sum_{k=1}^{n} k^6
  &= \tfrac17 n^7 + \tfrac12 n^6 + \tfrac12 n^5 - \tfrac16 n^3 + \tfrac1{42} n\\
\sum_{k=1}^{n} k^7
  &= \tfrac18 n^8 + \tfrac12 n^7 + \tfrac7{12} n^6 - \tfrac7{24} n^4
     + \tfrac1{12} n^2\\
\sum_{k=1}^{n} k^8
  &= \tfrac19 n^9 + \tfrac12 n^8 + \tfrac23 n^7 - \tfrac7{15} n^5
     + \tfrac29 n^3 - \tfrac1{30} n\\
\sum_{k=1}^{n} k^9
  &= \tfrac1{10} n^{10} + \tfrac12 n^9 + \tfrac34 n^8 - \tfrac7{10} n^6
     + \tfrac12 n^4 - \tfrac3{20} n^2\\
\sum_{k=1}^{n} k^{10}
  &= \tfrac1{11} n^{11} + \tfrac12 n^{10} + \tfrac56 n^9 - n^7 + n^5
     - \tfrac12 n^3 + \tfrac5{66} n.
\end{align*}
Notice that the coefficients in the second column are always $\tfrac12$, and
that after the third column the powers of $n$ with nonzero coefficients
decrease by~$2$ until $n^2$ or $n$ is reached. The coefficients in all but
the first two columns seem to be rather haphazard, but there actually is some
sort of pattern; finding it may be regarded as a super-perspicacity test. See
Problem~27-17 for the complete story.)

\pnum{8.} Prove that every natural number is either even or odd.

\pnum{9.} Prove that if a set $A$ of natural numbers contains $n_0$ and
contains $k+1$ whenever it contains $k$, then $A$ contains all natural
numbers $\geq n_0$.

\pnum{10.} Prove the principle of mathematical induction from the well-ordering
principle.

\pnum{11.} Prove the principle of complete induction from the ordinary
principle of induction. \textit{Hint:} If $A$ contains~$1$ and $A$ contains
$n+1$ whenever it contains $1, \dots, n$, consider the set $B$ of all $k$
such that $1, \dots, k$ are all in $A$.

\pnum{12.}
\begin{letters}
\item If $a$ is rational and $b$ is irrational, is $a+b$ necessarily
  irrational? What if $a$ and $b$ are both irrational?
\item If $a$ is rational and $b$ is irrational, is $ab$ necessarily
  irrational? (Careful!)
\item Is there a number $a$ such that $a^2$ is irrational, but $a^4$ is
  rational?
\item Are there two irrational numbers whose sum and product are both
  rational?
\end{letters}

\pnum{13.}
\begin{letters}
\item Prove that $\sqrt{3}$, $\sqrt{5}$, and $\sqrt{6}$ are irrational.
  \textit{Hint:} To treat $\sqrt{3}$, for example, use the fact that every
  integer is of the form $3n$ or $3n+1$ or $3n+2$. Why doesn't this proof
  work for $\sqrt{4}$?
\item Prove that $\sqrt[3]{2}$ and $\sqrt[3]{3}$ are irrational.
\end{letters}

\pnum{14.} Prove that
\begin{letters}
\item $\sqrt{2} + \sqrt{6}$ is irrational.
\item $\sqrt{2} + \sqrt{3}$ is irrational.
\end{letters}

\pnum{15.}
\begin{letters}
\item Prove that if $x = p + \sqrt{q}$ where $p$ and $q$ are rational, and
  $m$ is a natural number, then $x^m = a + b\sqrt{q}$ for some rational $a$
  and $b$.
\item Prove also that $(p - \sqrt{q})^m = a - b\sqrt{q}$.
\end{letters}

\pnum{16.}
\begin{letters}
\item Prove that if $m$ and $n$ are natural numbers and $m^2/n^2 < 2$, then
  $(m+2n)^2/(m+n)^2 > 2$; show, moreover, that
\[ \frac{(m+2n)^2}{(m+n)^2} - 2 < 2 - \frac{m^2}{n^2}. \]
\item Prove the same results with all inequality signs reversed.
\item Prove that if $m/n < \sqrt{2}$, then there is another rational number
  $m'/n'$ with $m/n < m'/n' < \sqrt{2}$.
\end{letters}

\pnum{*17.} It seems likely that $\sqrt{n}$ is irrational whenever the natural
number $n$ is not the square of another natural number. Although the method
of Problem~13 may actually be used to treat any particular case, it is not
clear in advance that it will always work, and a proof for the general case
requires some extra information. A natural number $p$ is called a
\textbf{prime number} if it is impossible to write $p = ab$ for natural
numbers $a$ and $b$ unless one of these is $p$, and the other~$1$; for
convenience we also agree that $1$ is \emph{not} a prime number. The first
few prime numbers are $2, 3, 5, 7, 11, 13, 17, 19$. If $n > 1$ is not a
prime, then $n = ab$, with $a$ and $b$ both $< n$; if either $a$ or $b$ is
not a prime it can be factored similarly; continuing in this way proves that
we can write $n$ as a product of primes. For example,
$28 = 4 \cdot 7 = 2 \cdot 2 \cdot 7$.

\begin{letters}
\item Turn this argument into a rigorous proof by complete induction. (To be
  sure, any reasonable mathematician would accept the informal argument, but
  this is partly because it would be obvious to her how to state it
  rigorously.)
\end{letters}

A fundamental theorem about integers, which we will not prove here, states
that this factorization is unique, except for the order of the factors. Thus,
for example, $28$ can never be written as a product of primes one of which is
$3$, nor can it be written in a way that involves $2$ only once (now you
should appreciate why $1$ is not allowed as a prime).

\begin{letters}[start=2]
\item Using this fact, prove that $\sqrt{n}$ is irrational unless $n = m^2$
  for some natural number $m$.
\item Prove more generally that $\sqrt[k]{n}$ is irrational unless $n = m^k$.
\item No discussion of prime numbers should fail to allude to Euclid's
  beautiful proof that there are infinitely many of them. Prove that there
  cannot be only finitely many prime numbers $p_1, p_2, p_3, \dots, p_n$ by
  considering $p_1 \cdot p_2 \cdots p_n + 1$.
\end{letters}

\pnum{*18.}
\begin{letters}
\item Prove that if $x$ satisfies
\[ x^n + a_{n-1} x^{n-1} + \cdots + a_0 = 0, \]
for some integers $a_{n-1}, \dots, a_0$, then $x$ is irrational unless $x$ is
an integer. (Why is this a generalization of Problem~17?)
\item Prove that $\sqrt{6} - \sqrt{2} - \sqrt{3}$ is irrational.
\item Prove that $\sqrt{2} + \sqrt[3]{2}$ is irrational. \textit{Hint:} Start
  by working out the first~$6$ powers of this number.
\end{letters}

\pnum{19.} Prove Bernoulli's inequality: If $h > -1$, then
\[ (1+h)^n \geq 1 + nh \]
for any natural number $n$. Why is this trivial if $h > 0$?

\pnum{20.} The Fibonacci sequence $a_1, a_2, a_3, \dots$ is defined as
follows:
\begin{align*}
a_1 &= 1,\\
a_2 &= 1,\\
a_n &= a_{n-1} + a_{n-2} \qquad \text{for } n \geq 3.
\end{align*}
This sequence, which begins $1, 1, 2, 3, 5, 8, \dots$, was discovered by
Fibonacci (circa 1175--1250), in connection with a problem about rabbits.
Fibonacci assumed that an initial pair of rabbits gave birth to one new pair
of rabbits per month, and that after two months each new pair behaved
similarly. The number $a_n$ of pairs born in the $n$th month is
$a_{n-1} + a_{n-2}$, because a pair of rabbits is born for each pair born the
previous month, and moreover each pair born two months ago now gives birth to
another pair. The number of interesting results about this sequence is truly
amazing---there is even a Fibonacci Association which publishes a journal,
\textit{The Fibonacci Quarterly}. Prove that
\[
a_n = \frac{
  \left(\dfrac{1+\sqrt{5}}{2}\right)^{n}
  - \left(\dfrac{1-\sqrt{5}}{2}\right)^{n}
}{\sqrt{5}}.
\]
One way of deriving this astonishing formula is presented in Problem~24-16.

\pnum{21.} The Schwarz inequality (Problem~1-19) actually has a more general
form:
\[
\sum_{i=1}^{n} x_i y_i
\leq
\sqrt{\sum_{i=1}^{n} x_i^2}\,
\sqrt{\sum_{i=1}^{n} y_i^2}.
\]
Give three proofs of this, analogous to the three proofs in Problem~1-19.

\pnum{22.} The result in Problem~1-7 has an important generalization: If
$a_1, \dots, a_n \geq 0$, then the ``arithmetic mean''
\[ A_n = \frac{a_1 + \cdots + a_n}{n} \]
and ``geometric mean''
\[ G_n = \sqrt[n]{a_1 \ldots a_n} \]
satisfy
\[ G_n \leq A_n. \]

\begin{letters}
\item Suppose that $a_1 < A_n$. Then some $a_i$ satisfies $a_i > A_n$; for
  convenience, say $a_2 > A_n$. Let $\bar{a}_1 = A_n$ and let
  $\bar{a}_2 = a_1 + a_2 - \bar{a}_1$. Show that
\[ \bar{a}_1 \bar{a}_2 \geq a_1 a_2. \]
Why does repeating this process enough times eventually prove that
$G_n \leq A_n$? (This is another place where it is a good exercise to provide
a formal proof by induction, as well as an informal reason.) When does
equality hold in the formula $G_n \leq A_n$?

The reasoning in this proof is related to another interesting proof.

\item Using the fact that $G_n \leq A_n$ when $n = 2$, prove, by induction on
  $k$, that $G_n \leq A_n$ for $n = 2^k$.
\item For a general $n$, let $2^m > n$. Apply part~(b) to the $2^m$ numbers
\[ a_1, \dots, a_n, \underbrace{A_n, \dots, A_n}_{2^m - n \text{ times}} \]
to prove that $G_n \leq A_n$.
\end{letters}

\pnum{23.} The following is a recursive definition of $a^n$:
\begin{align*}
a^1 &= a,\\
a^{n+1} &= a^n \cdot a.
\end{align*}
Prove, by induction, that
\begin{align*}
a^{n+m} &= a^n \cdot a^m,\\
(a^n)^m &= a^{nm}.
\end{align*}
(Don't try to be fancy: use either induction on $n$ or induction on $m$, not
both at once.)

\pnum{24.} Suppose we know properties P1 and P4 for the natural numbers, but
that multiplication has never been mentioned. Then the following can be used
as a recursive definition of multiplication:
\begin{align*}
1 \cdot b &= b,\\
(a+1) \cdot b &= a \cdot b + b.
\end{align*}
Prove the following (in the order suggested!):
\begin{align*}
a \cdot (b+c) &= a \cdot b + a \cdot c
  \quad\text{(use induction on $a$)},\\
a \cdot 1 &= a,\\
a \cdot b &= b \cdot a
  \quad\text{(you just finished proving the case $b = 1$)}.
\end{align*}

\pnum{25.} In this chapter we began with the natural numbers and gradually
built up to the real numbers. A completely rigorous discussion of this
process requires a little book in itself (see Part~V). No one has ever
figured out how to get to the real numbers without going through this
process, but if we do accept the real numbers as given, then the natural
numbers can be \emph{defined} as the real numbers of the form $1$, $1+1$,
$1+1+1$, etc. The whole point of this problem is to show that there is a
rigorous mathematical way of saying ``etc.''

\begin{letters}
\item A set $A$ of real numbers is called \textbf{inductive} if
\begin{steps}
\item $1$ is in $A$,
\item $k+1$ is in $A$ whenever $k$ is in $A$.
\end{steps}
Prove that
\begin{romans}
\item $\mathbf{R}$ is inductive.
\item The set of positive real numbers is inductive.
\item The set of positive real numbers unequal to $\tfrac12$ is inductive.
\item The set of positive real numbers unequal to~$5$ is not inductive.
\item If $A$ and $B$ are inductive, then the set $C$ of real numbers which
  are in both $A$ and $B$ is also inductive.
\end{romans}
\item A real number $n$ will be called a \textbf{natural number} if $n$ is in
  \emph{every} inductive set.
\begin{romans}
\item Prove that $1$ is a natural number.
\item Prove that $k+1$ is a natural number if $k$ is a natural number.
\end{romans}
\end{letters}

\pnum{26.} There is a puzzle consisting of three spindles, with $n$ concentric
rings of decreasing diameter stacked on the first (Figure~1). A ring at the
top of a stack may be moved from one spindle to another spindle, provided
that it is not placed on top of a smaller ring. For example, if the smallest
ring is moved to spindle~2 and the next-smallest ring is moved to spindle~3,
then the smallest ring may be moved to spindle~3 also, on top of the
next-smallest. Prove that the entire stack of $n$ rings can be moved onto
spindle~3 in $2^n - 1$ moves, and that this cannot be done in fewer than
$2^n - 1$ moves.

% Figure 1: Tower of Hanoi. Canonical SVG from
% assets/generators/calcu5-fig1.py; digits 2/3 are paths in the SVG, so
% Inkscape must not overlay LaTeX text. Width tracks \textwidth so the
% 243.6 mm native SVG fits the 130 mm digital / print-h measure.
\begin{center}
\includesvg[width=0.95\textwidth,inkscapelatex=false]{assets/rendered/calcu5-fig1}\\[0.6em]
\figlabel{1}
\end{center}

\pnum{*27.} University B. once boasted 17 tenured professors of mathematics.
Tradition prescribed that at their weekly luncheon meeting, faithfully
attended by all~17, any members who had discovered an error in their
published work should make an announcement of this fact, and promptly resign.
Such an announcement had never actually been made, because no professor was
aware of any errors in her or his work. This is not to say that no errors
existed, however. In fact, over the years, in the work of every member of the
department at least one error had been found, by some other member of the
department. This error had been mentioned to all other members of the
department, but the actual author of the error had been kept ignorant of the
fact, to forestall any resignations.

One fateful year, the department was augmented by a visitor from another
university, one Prof.~X, who had come with hopes of being offered a permanent
position at the end of the academic year. Naturally, he was apprised, by
various members of the department, of the published errors which had been
discovered. When the hoped-for appointment failed to materialize, Prof.~X
obtained his revenge at the last luncheon of the year. ``I have enjoyed my
visit here very much,'' he said, ``but I feel that there is one thing that I
have to tell you. At least one of you has published an incorrect result,
which has been discovered by others in the department.'' What happened the
next year?

\pnum{**28.} After figuring out, or looking up, the answer to Problem~27,
consider the following: Each member of the department already knew what
Prof.~X asserted, so how could his saying it change anything?
